\documentclass[11pt,fontset=fandol]{ctexart}
\usepackage[a4paper,margin=23mm]{geometry}
\usepackage{amsmath,amssymb,mathtools,bm,mathrsfs,esint}
\usepackage{graphicx,longtable,booktabs,array,calc}
\usepackage[normalem]{ulem}
\usepackage{xcolor,hyperref}
\hypersetup{hidelinks}
\setcounter{secnumdepth}{-1}
\newcounter{none}
\providecommand{\tightlist}{\setlength{\itemsep}{0pt}\setlength{\parskip}{0pt}}
\providecommand{\passthrough}[1]{#1}
\usepackage{amssymb,mathrsfs,esint,bm,newunicodechar}
\allowdisplaybreaks
\setlength{\emergencystretch}{3em}
\providecommand{\llangle}{\langle\!\langle}
\providecommand{\rrangle}{\rangle\!\rangle}
\newunicodechar{∞}{\ensuremath{\infty}}
\newunicodechar{∈}{\ensuremath{\in}}
\newunicodechar{∑}{\ensuremath{\sum}}
\newunicodechar{∫}{\ensuremath{\int}}
\newunicodechar{≤}{\ensuremath{\leq}}
\newunicodechar{≥}{\ensuremath{\geq}}
\newunicodechar{≠}{\ensuremath{\neq}}
\newunicodechar{→}{\ensuremath{\to}}
\newunicodechar{⇒}{\ensuremath{\Rightarrow}}
\newunicodechar{α}{\ensuremath{\alpha}}
\newunicodechar{β}{\ensuremath{\beta}}
\newunicodechar{γ}{\ensuremath{\gamma}}
\newunicodechar{λ}{\ensuremath{\lambda}}
\newunicodechar{φ}{\ensuremath{\varphi}}
\newunicodechar{ε}{\ensuremath{\varepsilon}}
\newunicodechar{σ}{\ensuremath{\sigma}}
\newunicodechar{θ}{\ensuremath{\theta}}
\newunicodechar{δ}{\ensuremath{\delta}}
\newunicodechar{η}{\ensuremath{\eta}}
\newunicodechar{π}{\ensuremath{\pi}}
\newunicodechar{ξ}{\ensuremath{\xi}}
\newunicodechar{Σ}{\ensuremath{\Sigma}}
\newunicodechar{Δ}{\ensuremath{\Delta}}
\newunicodechar{∂}{\ensuremath{\partial}}
\newunicodechar{∀}{\ensuremath{\forall}}
\newunicodechar{∃}{\ensuremath{\exists}}
\newunicodechar{√}{\ensuremath{\surd}}
\newunicodechar{−}{\ensuremath{-}}
\newunicodechar{μ}{\ensuremath{\mu}}
\newunicodechar{ρ}{\ensuremath{\rho}}
\newunicodechar{τ}{\ensuremath{\tau}}
\newunicodechar{ω}{\ensuremath{\omega}}
\newunicodechar{Ω}{\ensuremath{\Omega}}
\newunicodechar{χ}{\ensuremath{\chi}}
\newunicodechar{ψ}{\ensuremath{\psi}}
\newunicodechar{Γ}{\ensuremath{\Gamma}}
\newunicodechar{Φ}{\ensuremath{\Phi}}
\newunicodechar{Λ}{\ensuremath{\Lambda}}
\newunicodechar{ℝ}{\ensuremath{\mathbb{R}}}
\newunicodechar{ℤ}{\ensuremath{\mathbb{Z}}}
\newunicodechar{ℕ}{\ensuremath{\mathbb{N}}}
\newunicodechar{△}{\ensuremath{\triangle}}
\newunicodechar{⊗}{\ensuremath{\otimes}}
\begin{document}
\section{2012年杭州电子科技大学601数学分析考研真题}\label{ux5e74ux676dux5ddeux7535ux5b50ux79d1ux6280ux5927ux5b66601ux6570ux5b66ux5206ux6790ux8003ux7814ux771fux9898}

来源：\texttt{HDU/601数学分析/2012/2012年杭州电子科技大学601数学分析考研真题.pdf}。原
PDF 页序：1, 2。

\section{杭州电子科技大学}\label{ux676dux5ddeux7535ux5b50ux79d1ux6280ux5927ux5b66}

\section{2012年攻读硕士学位研究生入学考试}\label{ux5e74ux653bux8bfbux7855ux58ebux5b66ux4f4dux7814ux7a76ux751fux5165ux5b66ux8003ux8bd5}

\section{《数学分析》试题}\label{ux6570ux5b66ux5206ux6790ux8bd5ux9898}

\begin{quote}
助手整理：依据同题本地 OCR
恢复指数、分式及分段函数的漏项，保留题号与原题条件。
\end{quote}

（试题共四大题，2页，150分）

姓名 报考专业 准考证号

【所有答案必须写在答题纸上，做在试卷或草稿纸上无效！】

\subsection{一、填空题（每小题5分，共40分）}\label{ux4e00ux586bux7a7aux9898ux6bcfux5c0fux98985ux5206ux517140ux5206}

1.设
\(\varphi ( x ) = \int _ { x } ^ { x ^ { 3 } } \arctan { \frac { y } { x } } d y\)
，则φ'(x)=

\begin{enumerate}
\def\labelenumi{\arabic{enumi}.}
\setcounter{enumi}{1}
\tightlist
\item
  设 \(e^{xz}+xy=z^2\)，则 \(dz=\)\_\_\_\_。
\end{enumerate}

3．旋转抛物面 \(z = x ^ { 2 } + y ^ { 2 } - 1\)
在点(1,1,1)处切平面方程为

\begin{enumerate}
\def\labelenumi{\arabic{enumi}.}
\setcounter{enumi}{3}
\tightlist
\item
  \(\int _ { 1 } ^ { e } { \frac { d x } { x { \sqrt { 1 + \ln x } } } } =\)
\end{enumerate}

5.设 \(u = \arctan { \frac { x } { y } }\) ，则
\(\frac { { \partial } ^ { 2 } u } { { \partial } x ^ { 2 } }\)

6．交换积分次序
\(\int _ { 0 } ^ { 2 } d y \int _ { - \sqrt { 2 y - y ^ { 2 } } } ^ { \sqrt { 2 y - y ^ { 2 } } } f ( x , y ) d x\)
，则积分可表达为

7.函数
\(T ( s ) = \int _ { 0 } ^ { + \infty } x ^ { s - 3 } e ^ { - x } d x\)
的定义域为

8.设 \(y = \sin ^ { 2 } x\) ，则 \(y ^ { ( n ) } ( x ) =\)

\subsection{二、计算题（4小题，共30分）}\label{ux4e8cux8ba1ux7b97ux98984ux5c0fux9898ux517130ux5206}

\begin{enumerate}
\def\labelenumi{\arabic{enumi}.}
\tightlist
\item
  \(f ( x ) = x ^ { \sin x } + x { \sqrt { \sin { \sqrt { x } } } }\)
  ，求其一阶导数．（本小题7分）
\end{enumerate}

2．计算定积分
\(\int _ { - { \frac { \pi } { 4 } } } ^ { \frac { \pi } { 4 } } ( x ^ { 3 } \sin x ^ { 2 } + \cos ^ { 7 } 2 x )\)
dx．（本小题7分）

3．求函数 \(\int_0^x\dfrac{\cos t-1}{t}\,dt\)
的幂级数展开式并求其收敛域。（本小题8分）

4.计算
\(\int _ { L } ( x ^ { 2 } - y ) d x - ( x + \sin ^ { 2 } y ) d y\)
，其中L是在圆周 \(y = \sqrt { 2 x - x ^ { 2 } }\)
上由点(0,0)到(1,1)的一段弧.（本小题8分）

\subsection{三、综合题(每小题10分，共30
分）}\label{ux4e09ux7efcux5408ux9898ux6bcfux5c0fux989810ux5206ux517130-ux5206}

1．讨论函数y=x-2arctanx的单调区间、极值、凹凸区间、拐点及其渐近线.

2.设
\(f _ { n } ( x ) = x n ^ { k } e ^ { - n x } , x \in [ 0 , + \infty )\)
，确定k之范围，使函数列 \(\{ f _ { n } ( x ) \}\) 在
\([ 0 , + \infty ) -\) 致收敛.

3．讨论二元函数
\(f ( x , y ) = { \frac { x ^ { 2 } ( 1 + x ^ { 2 } ) - y ^ { 2 } ( 1 + y ^ { 2 } ) } { x ^ { 2 } + y ^ { 2 } } }\)
在原点处的二重极限和累次极限.

\subsection{四、证明题（每小题 10分，共 50
分)}\label{ux56dbux8bc1ux660eux9898ux6bcfux5c0fux9898-10ux5206ux5171-50-ux5206}

1.用ε-δ定义证明
\(\lim _ { x \to 3 } { \frac { x } { x ^ { 2 } + 1 } } = { \frac { 3 } { 1 0 } }\)

2．设 \(g\) 具有二阶连续导数，\(g(0)=0\)，证明函数
\(f(x)=\begin{cases}g(x)/x,&x\neq0,\\g'(0),&x=0\end{cases}\)
有连续的导函数。

3．证明：当x\textgreater0时，
\(\ln ( 1 + { \frac { 1 } { x } } ) > { \frac { 1 } { 1 + x } }\)

4．证明闭区间{[}a,b{]}上的单调函数f(x)必定黎曼可积.

5．证明：若在 \([ a , \infty ; c , d ]\) 内成立
\(\left| f ( x , y ) \right| \leq F ( x , y )\) ，并且关于
\(y \in [ c , d ]\) 积分\(\int _ { a } ^ { + \infty } F ( x , y ) d x\)
一致收敛，则 \(\int _ { a } ^ { + \infty } f ( x , y ) d x\) 关于
\(y \in [ c , d ]\) 亦一致收敛，且绝对收敛.
\end{document}
